%!TeX root = ../report.tex


\chapter{Anàlisi del Sistema}
\label{ch:analisi-del-sistema}

En aquest capítol es detalla l'anàlisi de requisits, restriccions i decisions de disseny que han guiat el
desenvolupament del projecte \textbf{Demochain}.
L'objectiu és oferir una visió comprensible de la infraestructura tant per a auditors externs com per a perfils tècnics
que no hagin participat directament en el procés de programació.
A continuació, es descriuen els actors del sistema, la recerca teòrica en criptografia i xarxes de consens en
contraposició a la implementació real de l'\ac{MVP}\@, les especificacions funcionals i no funcionals, les restriccions
imposades per l'entorn tecnològic i legal, i les decisions de disseny global de la plataforma.


\section{Identificació de Parts Interessades (\textit{Stakeholders})}\label{sec:stakeholders}

Per tal d'entendre l'abast i els objectius del sistema de votació, és imprescindible identificar quins són els actors
clau que intervenen en el procés electoral i quin benefici directe o indirecte obtenen de la seva implantació El sistema
s'ha concebut des d'una perspectiva modular on es contraposen els interessos dels administradors, electors, custodis i
verificadors per aconseguir un equilibri de seguretat.

\begin{enumerate}
    \item \textbf{Votants / electors que formen part del cens}:
    \begin{itemize}
        \item \textbf{Descripció}: Membres actius de la institució que convoca la votació, com ara els estudiants i
        professors en l'àmbit universitari.

        \item \textbf{Benefici}: Disposar d'una plataforma digital accessible que els permet exercir el sufragi de
        manera remota, evitant desplaçaments físics i esperes.
        El sistema els assegura la total privacitat de les seves preferències electorals i els proporciona mecanismes
        pràctics contra la coacció, com és el dret a modificar el vot emès mentre l'elecció estigui oberta.
        A més, en emprar una xarxa permissionada a cost zero per transacció, no han d'adquirir criptomonedes ni pagar
        taxes per votar.
        Finalment, tenen la capacitat de verificar individualment que la seva papereta s'ha registrat de forma íntegra a
        la xarxa d'Algorand mitjançant la interfície de consulta.
    \end{itemize}

    \item \textbf{Candidats i les diferents opcions de votació}:
    \begin{itemize}
        \item \textbf{Descripció}: Persones físiques que opten a un càrrec, llistes de partits o alternatives en un
        referèndum o consulta normativa.

        \item \textbf{Benefici}: Garantia que les regles del joc són immutables des de l'inici fins al final de la
        votació, ja que les dades s'emmagatzemen a una cadena de blocs (\textit{blockchain}), eliminant el risc de
        manipulació física d'urnes o pèrdua de paperetes.
        Així mateix, es beneficien d'un recompte automàtic i immediat un cop tancat el procés electoral, reduint el cost
        de campanya associat a la mobilització de delegats o interventors a les meses electorals tradicionals.
    \end{itemize}

    \item \textbf{Organitzadors (Comitè Electoral i Administradors)}:
    \begin{itemize}
        \item \textbf{Descripció}: Entitat o comitè electoral responsable de la gestió de les consultes, la qual cosa
        inclou crear l'organització, gestionar el cens d'adreces autoritzades, definir les opcions de les propostes i
        coordinar les dates d'inici i de tancament.

        \item \textbf{Benefici}: Reducció dràstica en els costos de gestió que acompanyen qualsevol sufragi tradicional
        (paperetes, urnes físiques, lloguer d'espais, dietes del personal de mesa, serveis de seguretat).
        L'automatització d'aquest procés elimina l'error humà i la necessitat d'escrutinis manuals llargs, oferint
        resultats transparents que reforcen la reputació democràtica de l'entitat.
    \end{itemize}

    \item \textbf{Custodis (\textit{Trustees})}:
    \begin{itemize}
        \item \textbf{Descripció}: Actors independents amb rols contraposats dins de la comunitat (com ara membres de la
        junta de govern, representants dels estudiants o auditors externs) que reben un fragment de la clau de
        desxifratge del recompte global.

        \item \textbf{Benefici}: Control conjunt de l'escrutini final.
        Cap entitat per si sola pot desxifrar els vots individuals ni calcular els resultats abans d'hora.
        Gràcies a un esquema de llindar, poden participar de manera asíncrona per reconstruir el resultat, obtenint
        proves matemàtiques que validen la seva honestedat i actuació davant els electors sense por de ser acusats de
        manipulació.
    \end{itemize}

    \item \textbf{Auditors i observadors independents}:
    \begin{itemize}
        \item \textbf{Descripció}: Entitats externes, periodistes o observadors internacionals que vetllen per la
        legalitat i legitimitat del sufragi.

        \item \textbf{Benefici}: Accés directe i en temps real a tot l'historial de transaccions d'Algorand i al segell
        d'Ethereum Sepolia sense necessitar permisos explícits ni demanar autorització als administradors.
        Això els permet auditar de cap a fi la integritat del procés reexecutant l'algorisme de Schulze directament
        sobre les caixes de dades públicament disponibles.
    \end{itemize}

    \item \textbf{Entitats Democràtiques i DAO (Adoptants B2B)}:
    \begin{itemize}
        \item \textbf{Descripció}: Sindicats, associacions, cooperatives o organitzacions descentralitzades que cerquen
        una eina de govern electrònic segura.

        \item \textbf{Benefici}: Reutilització d'aquesta tecnologia de codi obert per a organitzar els seus processos de
        presa de decisions interns a cost zero de desenvolupament, adaptant el disseny modular a les seves estructures
        de govern.
    \end{itemize}

    \item \textbf{Equip de desenvolupament i comunitat acadèmica}:
    \begin{itemize}
        \item \textbf{Descripció}: Investigadors i programadors que dissenyen la solució i avaluen la viabilitat de
        noves primitives de xifratge o analitzen el comportament de l'algorisme de Schulze \textit{on-chain}.

        \item \textbf{Benefici}: Disposen d'un cas pràctic i d'un repositori de codi modular escrit en Python i
        Solidity.
        Per a ells, aquest projecte ha constituït un gran entorn de formació pràctica, atès que no tenien pràcticament
        cap coneixement previ sobre tecnologia \textit{blockchain} o criptografia abans de començar les catorze setmanes
        de desenvolupament.
    \end{itemize}
\end{enumerate}


\section{Recerca Criptogràfica i de Consens per a un Disseny Robust}
\label{sec:recerca-criptografica}

El disseny de qualsevol sistema de votació electrònica s'enfronta a un dilema fonamental: com garantir que només els
electors autoritzats votin i que cadascú ho faci una sola vegada (integritat i autenticació), alhora que s'assegura que
ningú, ni tan sols l'administrador del sistema, pugui vincular un vot amb la identitat de qui l'ha emès (privacitat),
tot permetent que qualsevol ciutadà comprovi de forma independent que el recompte és correcte (auditabilitat pública).
En els sistemes de votació tradicionals, aquest equilibri es confia a mesures físiques (urnes segellades, sobres opacs i
meses electorals multipartidistes).
Traslladar aquestes garanties al món digital sense dependre d'una autoritat central que podria ser corrupta requereix
l'ús de primitives criptogràfiques avançades.

Durant la fase inicial del projecte, vàrem dur a terme una recerca profunda sobre l'estat de l'art en protocols de
votació electrònica segura.
En aquesta recerca, el projecte de codi obert \textbf{DaVinci} (disponible a \url{https://davinci.vote/}) ens va servir
de guia i referència per a estructurar i entendre com coordinar diferents tecnologies de consens i privacitat.
A partir d'aquest referent, vàrem identificar els conceptes criptogràfics clau que haurien de formar part d'una
arquitectura teòrica robusta:

\begin{itemize}
    \item \textbf{Proveïdor d'Identitat (IdP - Identity Provider)}: Representa l'entitat de confiança que gestiona el
    \item cens real dels electors (per exemple, el sistema d'autenticació central d'una universitat).
    En un disseny plenament descentralitzat, l'\ac{IdP} signa una credencial criptogràfica per a cada usuari registrat que
    demostra el seu dret al sufragi, sense vincular-lo directament a una adreça de bloc pública (\textit{wallet}).

    \item \textbf{Proves de Coneixement Zero (ZK-SNARK)}: Són proves criptogràfiques que permeten a un elector provar
    davant el contracte intel·ligent que té una credencial vàlida signada per l'\ac{IdP} (i, per tant, forma part del
    cens) sense revelar la seva identitat real ni revelar la signatura de l'\ac{IdP}\@.
    A nivell conceptual, el proveïdor executa un circuit amb entrades privades (les seves credencials) i paràmetres
    públics per a generar una prova compacta que qualsevol node pot verificar de forma instantània.
    Per evitar el doble vot sense comprometre l'anonimat, es va considerar l'ús d'un \textit{nullifier} (anul·lador), un
    identificador determinista derivat de la clau privada del votant i de la proposta.
    El contracte intel·ligent rebutja transaccions amb anul·ladors ja utilitzats.
    No obstant això, el tema del \textit{nullifier} no el vàrem acabar de concretar bé en la nostra recerca; dissenyar
    una estructura robusta i segura per a aquesta primitiva és extremadament complex i no vàrem trobar un disseny
    definitiu que s'adaptés de forma senzilla a la nostra infraestructura.

    \item \textbf{Xifratge Homomòrfic}: És un tipus de xifratge que permet efectuar operacions matemàtiques directament
    sobre dades xifrades (\textit{ciphertexts}) de manera que, en desxifrar el resultat, s'obté la suma o l'operació en
    clar de les dades de partida.
    En el context electoral, les preferències de cada votant es xifren en el dispositiu client abans de ser enviades.
    La xarxa de blocs agrega de forma multiplicativa totes les paperetes xifrades directament \textit{on-chain},
    generant un únic resultat xifrat de toda la votació.
    Això assegura que les paperetes individuals mai es desxifren, protegint de forma absoluta el secret del sufragi.

    \item \textbf{Custodis o Trustees}: Són entitats de confiança encarregades de la custòdia de la clau de desxifratge
    global.
    Per evitar un únic punt de fallada o que l'administrador del sistema pugui desxifrar els vots abans de temps,
    s'empra un esquema de llindar.
    Mitjançant un protocol de generació de claus distribuïdes (DKG), la clau privada de desxifratge es divideix en
    fragments (\textit{shares}) repartits entre diversos custodis.
    Per desxifrar el recompte homomòrfic acumulat un cop tancat el procés, es requereix que un nombre mínim o llindar de
    custodis (com a manco la majoria absoluta) aporti la seva \textit{share} parcial.
\end{itemize}

\subsection{Desviació del Disseny i Limitacions del Prototip MVP}
\label{subsec:desviacio-mvp}

Malgrat haver definit aquest disseny teòric, l'equip va haver de fer front a una important limitació pràctica: la manca
de coneixements previs profunds sobre criptografia avançada i la dificultat extrema que suposa implementar i auditar tot
aquest flux de forma totalment segura.
Dissenyar circuits de coneixement zero, coordinar la cerimònia de DKG entre custodis i desplegar operacions
homomòrfiques a la màquina virtual de forma eficient no és gens trivial i requereix mesos de desenvolupament
especialitzat.

Per mor d'això, es va decidir orientar el desenvolupament del prototip \ac{MVP} cap a una arquitectura simplificada que
fos viable dins els terminis del projecte.
Aquesta desviació de disseny implica els següents canvis:

\begin{enumerate}
    \item \textbf{Vots en clar}: Les paperetes de preferències ordinals es registren i s'emmagatzemen directament en
    clar a les \textit{boxes} d'Algorand, sense xifratge homomòrfic.
    Qualsevol pot auditar els vots reals consultant la \textit{blockchain}.

    \item \textbf{Cens identificat}: La gestió d'electors s'efectua a través d'una llista blanca d'adreces de carteres
    físiques (\textit{wallets}) gestionada per l'organitzador, prescindint de les proves \ac{ZK-SNARK} i de
    l'anul·lador.
    Els usuaris s'identifiquen i signen els seus vots amb la seva pròpia clau pública d'Algorand.

    \item \textbf{Càlcul fora de la cadena}: Ateses les limitacions de \textit{bytecode} d'Algorand, el recompte i
    l'algorisme de Schulze es realitzen al navegador (\textit{client-side}) obtenint els vots directament de les
    \textit{boxes} del contracte.
    La cerimònia de desxifratge de custodis s'ha posposat completament com a treball futur.
\end{enumerate}

Aquesta arquitectura de l'\ac{MVP} permet disposar d'un sistema funcional que garanteix la immutabilitat física dels vots
emesos a Algorand i l'ancoratge del resultat final a Ethereum Sepolia, al preu d'ajornar el secret criptogràfic del vot
i l'anònim absolut de l'elector per a futures iteracions de producció de la plataforma.


\section{Especificació Funcional}
\label{sec:especificacio-funcional}

El comportament del sistema es defineix a partir de les interaccions entre els diferents rols de la plataforma,
detallant un flux de travail continu que guia tot el procés.

\subsection{Cicle de Vida del Sufragi i Casos d'Ús}
\label{subsec:casos-d-us}

El cicle de vida d'una consulta electoral a la plataforma es defineix a partir de les interaccions dels diferents actors
i del flux lògic del sistema, el qual es pot esquematitzar en els següents vuit casos d'ús principals:

\begin{itemize}
    \item \textbf{UC-1: Crear Organització}: Qualsevol usuari registrat que connecti la seva cartera de criptomonedes
    (\textit{wallet}) pot registrar una nova organització a la xarxa d'Algorand proporcionant un nom i una descripció,
    esdevenint automàticament l'organitzador i propietari del registre.

    \item \textbf{UC-2: Gestionar Cens}: L'organitzador administra els electors autoritzats afegint o eliminant adreces
    de \textit{wallet} de la llista blanca de l'entitat.
    La interfície web gestiona aquesta càrrega en lots de com a màxim set adreces per transacció per evitar un consum
    excessiu de memòria de la màquina virtual.

    \item \textbf{UC-3: Crear Proposta}: Qualsevol membre registrat al cens de l'organització pot iniciar una consulta
    detallant el títol, la descripció, la llista de candidats ordinals (que ha de ser com a mínim de dos), la data
    d'inici i la data de finalització del sufragi.

    \item \textbf{UC-4: Emetre Vot d'Aprovació}: Durant una fase preliminar, els electors aproven o rebutgen la
    viabilitat de la proposta.
    Aquest vot d'aprovació és immutable i irrevocable, i només si s'assoleix un quòrum de dues terceres parts ($2/3$) de
    vots favorables abans de la data d'inici, la proposta passa a la fase electoral.

    \item \textbf{UC-5: Emetre/Modificar Vot Preferencial}: Un cop s'arriba a la data d'inici i durant la finestra
    activa, els votants ordenen els candidats per preferències mitjançant una interfície d'arrossegar i deixar anar
    (\textit{drag \& drop}).
    Es permet enviar noves paperetes per sobreescriure el registre anterior a les \textit{boxes} del contracte i així
    evitar coaccions.

    \item \textbf{UC-6: Recompte i Escrutini}: En tancar-se la votació, es realitza el recompte.
    Conforme al disseny simplificat de l'\ac{MVP}, les paperetes es llegeixen directament de les \textit{boxes}
    d'Algorand en clar i l'escrutini es resol mitjançant l'aplicació client que executa l'algorisme de Schulze.

    \item \textbf{UC-7: Ancoratge postelectoral}: Un servei independent recopila les paperetes finals d'Algorand, les
    ordena lexicogràficament per adreça, genera un \ac{JSON} canònic sense espais, calcula el \textit{hash}
    \ac{SHA}-256 de 32 bytes i l'envia al contracte notari d'Ethereum Sepolia per certificar el recompte.

    \item \textbf{UC-8: Auditabilitat i Verificació Pública}: Qualsevol ciutadà o observador extern pot consultar de
    manera individual la seva papereta a Algorand i comprovar si el \textit{hash} del resultat global coincideix amb el
    \textit{hash} d'ancoratge validat pel consens de les universitats a Ethereum Sepolia.
\end{itemize}

A continuació, es mostra el diagrama de flux que representa gràficament la interacció dels diferents rols amb cada cas
d'ús:

% TODO

\subsection{Metodologia del Backlog i Històries d'Usuari}
\label{subsec:user-stories}

El disseny de la lògica funcional i dels rols descrits s'ha planificat mitjançant metodologies àgils utilitzant GitHub
CLI, connectant de manera directa amb el repositori de codi.
La jerarquia de treball s'ha estructurat en èpiques, històries d'usuari, \textit{spikes} de recerca, tasques
específiques i gestió d'errors o \textit{bugs}, cosa que ha facilitat la distribució de feines durant les catorze
setmanes de durada del projecte.

A continuació es detallen les històries d'usuari clau que s'han emprat per a guiar el desenvolupament:

\begin{itemize}
    \item \textbf{US-1 (Gestió d'organitzacions i cens)}: Com a organitzador, vull poder crear una organització i
    gestionar el cens d'electors mitjançant càrregues en lot o fitxers \ac{CSV}, per tal de preparar les votacions de
    manera àgil i sense errors.

    \item \textbf{US-2 (Autenticació i connexió)}: Com a elector, vull poder vincular la meva cartera Pera Wallet de
    forma transparent des del navegador, per tal d'autenticar-me i signar les transaccions de votació de forma nativa i
    senzilla.

    \item \textbf{US-3 (Interfície de vot preferencial)}: Com a votant, vull disposar d'una interfície d'arrossegar i
    deixar anar (\textit{drag \& drop}) per tal d'ordenar els candidats segons les meves preferències, fent que el meu
    vot sigui més expressiu que en un sistema majoritari simple.

    \item \textbf{US-4 (Protecció contra la coacció)}: Com a elector, vull poder efectuar nous enviaments que
    sobreescriguin totalment la meva papereta anterior a les \textit{boxes} del contracte de manera opaca per a un
    observador extern, per tal de protegir-me contra la coacció o la compravenda de vots.

    \item \textbf{US-5 (Seguretat de la custòdia)}: Com a trustee o custodi, vull poder registrar-me en el sistema per a
    custodiar el procés, entenent que en futures versions de la plataforma podré aportar proves de coneixement zero i
    shares de desxifratge.

    \item \textbf{US-6 (Auditabilitat de l'ancoratge)}: Com a auditor acadèmic, vull poder verificar directament a
    Ethereum Sepolia si un resultat de votació té la certificació de consens del consorci d'universitats, per tal de
    validar formalment la legitimitat del sufragi.
\end{itemize}

Els \textit{spikes} de recerca inicials es van utilitzar per estudiar la viabilitat de dissenyar un cens totalment
anònim \textit{on-chain} i coordinar de forma eficient la cerimònia de desxifratge dels custodis abans de desglossar les
tasques de desenvolupament i programar els diferents components de la plataforma.

\subsection{Requisits de Producte}
\label{subsec:requisits-producte}

Els requisits de producte del sistema asseguren la consistència jurídica i tècnica del sufragi:

\begin{itemize}
    \item \textbf{RP-1 (Marge temporal de creació)}: La data d'inici d'una proposta ha d'estar, com a manco, a tres dies
    de distància de la seva creació en el contracte intel·ligent.

    \item \textbf{RP-2 (Finestra electoral mínima)}: La durada activa del procés de votació preferencial ha de tenir una
    durada d'almenys vint-i-quatre hores per garantir la participació.

    \item \textbf{RP-3 (Cens de dades atòmic)}: El registre del cens ha de validar que no s'introdueixen adreces
    duplicades ni adreces nul·les (adreça zero) a la llista blanca de l'organització.

    \item \textbf{RP-4 (Consistència de la papereta)}: El contracte intel·ligent d'Algorand ha de rebutjar de forma
    estricta paperetes de preferències que no continguin exactament el mateix nombre d'elements que candidats té la
    proposta.

    \item \textbf{RP-5 (Estat terminal de l'urna)}: Un cop calculat el guanyador i tancat l'escrutini de la proposta,
    l'elecció canvia al seu estat immutable terminal, tancant de forma completa qualsevol escriptura posterior.
\end{itemize}


\section{Especificació No Funcional i Paràmetres de Qualitat}
\label{sec:especificacio-no-funcional}

Els requisits no funcionals del sistema de votació estan orientats a garantir la seguretat, la disponibilitat i
l'auditabilitat de les consultes, alhora que busquen optimitzar l'eficiència computacional.

\subsection{Requisits No Funcionals}\label{subsec:requisits-no-funcionals}

\begin{itemize}
    \item \textbf{RNF-1 (Seguretat i secret del sufragi)}: Els vots individuals s'han de protegir per a mantenir-ne
    l'anonimat i el secret de les paperetes a la xarxa d'Algorand mitjançant adreces de cartera pseudònimes.

    \item \textbf{RNF-2 (Immutabilitat)}: Qualsevol transacció confirmada (vot emès, canvi de cens, creació
    d'organització) esdevé completament inalterable gràcies al protocol de consens immediat d'Algorand, impedint
    qualsevol modificació posterior.

    \item \textbf{RNF-3 (Disponibilitat)}: L'urna digital ha de ser accessible de forma permanent i ininterrompuda
    (24/7) durant tota la finestra de votació, independentment d'eventuals caigudes dels servidors de la institució
    convocant.

    \item \textbf{RNF-4 (Auditabilitat de cap a fi)}: Es fomenta l'auditabilitat \textit{end-to-end}, de manera que
    qualsevol observador o elector pot descarregar i verificar de forma asíncrona que la seva papereta s'ha
    comptabilitzat correctament i que el \textit{hash} coincideix amb el segell d'Ethereum Sepolia.

    \item \textbf{RNF-5 (Eficiència Computacional i cost)}: El disseny de l'execució s'optimitza per evitar un consum
    elevat d'operacions \textit{on-chain} a Algorand i redueix l'escriptura a Ethereum a una única transacció global per
    elecció finalitzada per evitar uns costos de gas inacceptables.
\end{itemize}

\subsection{Paràmetres de Qualitat}
\label{subsec:parametres-qualitat}

Per mesurar la qualitat i la resiliència del sistema de seguretat distribuïda s'apliquen tres paràmetres clau:

\begin{itemize}
    \item \textbf{Llindar de Seguretat de Custòdia ($K$-de-$N$)}: Per poder desxifrar la matriu agregada final de
    preferències sota l'esquema de seguretat de llindar del disseny teòric, cal la participació de com a mínim la
    majoria absoluta dels custodis autoritzats, la qual cosa és causada per la relació:
    \[ K = \left\lceil \frac{N + 1}{2} \right\rceil \]
    Aquesta equació assegura que fa falta el consens de la majority absoluta dels custodis per desbloquejar els
    resultats de la votació, evitant que una minoria o un únic administrador pugui comprometre el procés de forma
    unilateral.

    \item \textbf{Consens d'Ancoratge a Ethereum}: El contracte notari de Sepolia exigeix que el \textit{hash} enviat
    del resultat final rebi el consens de dues terceres parts ($2/3$) dels nodes universitaris autoritzats abans de
    marcar l'elecció com a oficialment ancorada.
    El llindar $K_{notary}$ es congela mitjançant la fórmula:
    \[ K_{notary} = \left\lceil \frac{2N + 2}{3} \right\rceil \]

    \item \textbf{Temps de Confirmació}: Aprofitant les característiques tècniques de la xarxa permissionada d'Algorand,
    qualsevol transacció de vot o d'aprovació s'ha de veure confirmada en un temps màxim de quatre segons, evitant
    retards en l'escriptura de dades i sense el risc de patir reorganitzacions probabilístiques de blocs.
\end{itemize}


\section{Restriccions}
\label{sec:restriccions}

El disseny arquitectònic de Demochain ha hagut d'adaptar-se a una sèrie de limitacions imposades pel \textit{software} de les
xarxes utilitzades, els reglaments legals de protecció de dades i el context d'aprenentatge de l'equip.

\subsection{Restriccions Tecnològiques}
\label{subsec:restriccions-tecnologiques}

El disseny de la lògica \textit{on-chain} s'ha enfrontat a límits físics estrictes imposats per la Màquina Virtual
d'Algorand (\ac{AVM}) i el protocol d'Ethereum.
Inicialment, l'\ac{AVM} limita la mida màxima del \textit{bytecode} dels contractes intel·ligents a només 8KB de
memòria.
Aquesta severa restricció de mida impedeix codificar directament \textit{on-chain} algorismes complexos de cerca de
camins ponderats en grafs de preferències, com ara l'algorisme de Floyd-Warshall que fa servir el mètode Schulze, el
qual té una complexitat computacional de triple bucle de tipus $O(n^3)$.
Per mor d'aquest límit de \textit{bytecode}, el sistema s'ha vist obligat a delegar la resolució del guanyador a
l'aplicació client (fora de la cadena de blocs), de manera que el navegador llegeix les paperetes, calcula el mètode
Schulze i les \textit{boxes} del contracte intel·ligent es limiten a persistir les dades de forma estructurada.

En segon lloc, l'emmagatzematge en \textit{boxes} d'Algorand implica un cost anomenat Saldo Mínim Requerit (MBR) per
byte per evitar la saturació del registre d'estat de la xarxa.
Aquest cost en microAlgos que el mateix contracte ha de mantenir com a fons de reserva es calcula mitjançant la fórmula:
\[ \text{MBR (microAlgos)} = 2500 + 400 \times (\text{len(clau)} + \text{mida\_box}) \]
Com que cada organització, membre del cens, proposta i vot emès genera noves caixes de dades al contracte, cal gestionar
l'MBR com a paràmetre crític per no deixar el compte del contracte sense saldos suficients.
A més, hi ha una restricció que impedeix que una transacció faci referència a més de vuit caixes simultàniament, la qual
cosa condiciona la manera en què el client ha d'enviar les paperetes a la xarxa.
Finalment, el cost econòmic extremadament alt de fer crides a Ethereum L1 obliga a utilitzar la xarxa de proves Sepolia
per a les universitats, i el protocol només envia un únic \textit{hash} de 32 bytes (\texttt{bytes32}) per cada elecció
un cop s'ha completat tot l'escrutini.

Sense segones capes o agregació de transaccions (\textit{rollups}), Algorand té un límit físic d'unes 6.000 transaccions
per segon.
És una xifra molt bona per a entitats mitjanes com universitats o empreses, però insuficient per a gestionar eleccions
estatals simultànies on participen milions de ciutadans.

\subsection{Restriccions Legals i de Privacitat}
\label{subsec:restriccions-legals}

En l'àmbit legal, el disseny s'enfronta directament a la normativa europea de protecció de dades (RGPD).
Atès que la cadena de blocs és històricament immutable per disseny, és tècnicament impossible esborrar les dades un cop
confirmades a la xarxa, fet que s'oposa al dret a l'oblit de la normativa legal.
Per resoldre aquesta restricció, el sistema prohibeix de manera absoluta guardar \textit{on-chain} qualsevol tipus de
dada personal identificable, com ara noms, correus o documents nacionals d'identitat.
La identificació física del votant es realitza exclusivament al Proveïdor d'Identitat (\ac{IdP}) centralitzat de la
institució convocant.

D'altra banda, el secret del vot s'oposa de forma frontal a la total transparència d'una \textit{blockchain} on
l'historial de transaccions de qualsevol adreça és pública per a tothom.
Com s'ha explicat a la secció de recerca criptogràfica, la privacitat teòrica mitjançant \ac{ZK-SNARK} i xifratge
homomòrfic ha hagut de ser simplificada en l'\ac{MVP} actual, on l'anonimat dels electors es recolza únicament en l'ús
d'adreces de cartera genèriques i pseudònimes, esdevenint una limitació de seguretat assumida pel prototip.

\subsection{Restriccions d'Aprenentatge i Terminis del Projecte}
\label{subsec:restriccions-aprenentatge}

La principal restricció que ha condicionat el disseny i el calendari del projecte ha estat la manca total d'experiència
prèvia en tecnologia \textit{blockchain} i criptografia per part de l'equip de desenvolupament abans de començar les
catorze setmanes de termini establertes:

\begin{itemize}
    \item Es va haver de dur a terme una intensa recerca inicial (\textit{spikes} de recerca) des de zero per copsar
    conceptes de distribució de dades, signatura de transaccions, gestió de carteres i funcionament de les màquines
    virtuals (\ac{AVM} i \ac{EVM}).

    \item Aquest escenari d'aprenentatge intensiu va obligar a prendre decisions realistes.
    Com s'ha detallat a la Secció~\ref{sec:recerca-criptografica}, el desenvolupament segur de circuits \ac{ZK-SNARK},
    el model d'anul·lador i el protocol de desxifratge homomòrfic amb DKG de custodis requereix competències altament
    especialitzades i un temps de programació inexistent en el calendari.
    Per consegüent, per assegurar el lliurament d'un prototip funcional de cap a fi (\ac{MVP}), es va decidir ajornar la
    implementació criptogràfica com a feina futura, centrant els esforços a estabilitzar els fluxos de cens, votació
    preferencial i ancoratge inter-blockchain.
\end{itemize}


\section{Decisions de Disseny Global}
\label{sec:decisions-disseny-global}

A l'hora de conceptualitzar Demochain, s'han pres decisions clau sobre l'arquitectura global que prioritzen la
transparència i la resistència a la manipulació per sobre de models centralitzats convencionals.
La decisió d'utilitzar Algorand en lloc d'Ethereum o altres plataformes de primera capa es va basar en criteris
econòmics i de facilitat de desenvolupament:

\begin{enumerate}
    \item \textbf{Comissions extremadament econòmiques}: Les comissions d'Algorand són de només 0,001 ALGO per
    transacció, la qual cosa feia factible el projecte fins i tot amb xarxes públiques.

    \item \textbf{Eliminació de barreres d'entrada gràcies a Python}: Poder programar els contractes intel·ligents
    directament en Python natiu mitjançant el compilador Puyapy (\texttt{algopy}) va permetre a un equip que no coneixia
    la \textit{blockchain} començar a desenvolupar des del primer dia, estalviant-se l'aprenentatge de llenguatges
    complexos com Solidity o Rust.
\end{enumerate}

Durant el cicle de desenvolupament de Demochain, el model de xarxa ha patit una evolució molt interessant arran de la
identificació de noves limitacions.
Inicialment, es va plantejar un model on els mateixos telèfons mòbils dels electors actuessin com a nodes validators de
la xarxa.
Es pensava que mitjançant el mecanisme de VRF (Funció Aleatòria Verificable) d'Algorand es podria obligar cada votant a
mantenir l'aplicació activa almenys dos minuts per poder votar, obligant-los a participar del consens de forma efímera.
L'estudi tècnic posterior va demostrar que això era completament inviable a la pràctica per mor de la inestabilitat de
les xarxes mòbils, el desgast de bateria dels terminals i els exigents requisits de bloqueig de fons (\textit{staking})
que té el protocol de consens de la xarxa.

Davant d'aquesta impossibilitat, es va dissenyar una infraestructura de xarxa privada permissionada Els nodes validators
passaven a ser gestionats per un consorci d'entitats de confiança recíproca amb interessos contraposats (per exemple,
les diferents universitats de la comunitat autònoma).
Això permetia un control de participació i fixar les comissions de la xarxa a zero de manera nativa (\texttt{min\_fee =
0}), eliminant la necessitat de finançar les carteres dels votants.

Més endavant, es va intentar adaptar el \textit{software} per a funcionar directament sobre la xarxa pública d'Algorand
(testnet/mainnet).
Ens vàrem adonar que la xarxa pública oferia grans avantatges de \textit{liveness}, seguretat contra atacs de denegació
de consens locals i una transparència absoluta per a auditors mundials.
No obstant això, es va topar amb un obstacle: com que les comissions són obligatòries a la xarxa pública (malgrat que
siguin negligibles), l'única forma d'evitar que el votant hagués d'adquirir criptomonedes era crear un servei
centralitzat de patrocini (\textit{sponsor/relayer}) que cobrís el cost de les transaccions a través d'un grup atòmic.
Això trencava la naturalesa descentralitzada de la votació (el patrocinador centralitzat podia censurar vots lliurement
no signant-ne el patrocini).
A més a més, l'equip ja tenia enllestit i testejat tot el mòdul d'ancoratge (\textit{anchoring}) a Ethereum Sepolia,
dissenyat precisament per a compensar la menor resiliència de seguretat d'una xarxa privada.

Finalment, es va decidir mantenir el disseny de xarxa privada permissionada controlada pel consorci d'universitats.
Tanmateix, es va decidir conservar la integració amb \textbf{Pera Wallet} a la interfície client.
Aquesta decisió de disseny manté viva la possibilitat de publicar l'aplicació en un futur a la mainnet pública
d'Algorand, ja que el \textit{software} és plenament compatible amb carteres externes; només caldria implementar un
mètode de distribució inicial de saldos o acceptar un relayer distribuït de comissions.


\section{Mètode de Recompte}
\label{sec:metode-recompte}

Una de les eleccions clau en el disseny de Demochain és el mètode per avaluar les preferències ordinals i determinar-ne
el guanyador de la consulta.

\subsection{El Teorema d'Impossibilitat d'Arrow}
\label{subsec:teorema-arrow}

Dins de la teoria de l'elecció social, el \textbf{Teorema de la Impossibilitat d'Arrow} demostra matemàticament que cap
sistema de votació ordinal pot complir simultàniament un conjunt de cinc criteris o condicions de justícia electoral
quan hi ha tres o més candidats:

\begin{enumerate}
    \item \textbf{Domini no restringit (Universalitat)}: El sistema ha d'admetre qualsevol perfil o combinació possible
    de preferències ordinals que els electors decideixin introduir.
    Els votants tenen llibertat absoluta per a ordenar els candidats com desitgin.

    \item \textbf{Absència de dictador (No dictadura)}: No hi ha d'haver cap elector que determini de forma unilateral
    el resultat del sistema independentment de les preferències de la resta de votants.
    El comportament col·lectiu ha de reflectir múltiples veus.

    \item \textbf{Eficiència de Pareto (Unanimitat)}: Si absolutament tots els electors prefereixen de forma estricta el
    candidat X per sobre del candidat Y, l'ordre col·lectiu del recompte ha de situar X abans que Y\@.

    \item \textbf{Independència d'alternatives irrellevants (IIA)}: La posició de preferència relativa entre dos
    candidats X i Y en el recompte final només ha de dependre de com els electors valorin X i Y entre ells, sense que
    s'alteri per la presència, introducció o eliminació d'un tercer candidat Z\@.

    \item \textbf{Monotonicitat}: Si un elector decideix modificar la seva papereta per a elevar la posició o el suport
    a una candidatura concreta, la posició col·lectiva final d'aquesta candidatura no ha de decaure de cap de les
    maneres (és a dir, donar més suport no ha de perjudicar mai el candidat).
\end{enumerate}

Aquest teorema demostra matemàticament que \textbf{no existeix un sistema electoral ordinal perfecte}, i que qualsevol
decisió de disseny representa una tria de compromís (\textit{trade-off}) on s'ha de sacrificar qualque criteri de
justícia teòrica per a obtenir viabilitat pràctica.

\subsection{Comparativa dels Mètodes de Votació i el Criteri de Condorcet}
\label{subsec:comparativa-metodes}

El \textbf{criteri de Condorcet} és un principi clau en la teoria d'elecció social.
Estipula que si un candidat és preferit de forma majoritària en tots els enfrontaments directes o per parelles (duels en
un-contra-un) davant de cadascun dels altres candidats, aquest candidat s'anomena \textbf{guanyador de Condorcet}.
En conseqüència, qualsevol mètode electoral que es defineixi com a just hauria de seleccionar aquest guanyador com a
vencedor de l'elecció.

Per justificar l'ús del \textbf{mètode Schulze}, vàrem fer una comparació detallada amb altres sistemes comuns:
\begin{itemize}
    \item \textbf{Sistema majoritari simple (vot únic)}: Molt inexpressiu.
    L'elector només pot triar un candidat, perdent tota informació sobre les seves altres preferències.
    Sovint força els votants al vot estratègic i no sol triar el guanyador de Condorcet.

    \item \textbf{Mètode de Borda}: Tot i ser ordinal i senzill (es donen punts en funció de la posició), viola el
    criteri de Condorcet.
    A més, és extremadament vulnerable a atacs de manipulació de clons, on s'introdueixen candidats pràcticament
    idèntics per a inflar artificialment la puntuació d'una opció i distorsionar el resultat.

    \item \textbf{Vot Transferible o Desfilament Instantani (IRV)}: S'elimina el candidat amb manco vots de primera
    preferència i es transfereixen els vots fins a trobar majoria.
    El defecte clau és que no és monòton (viola el criteri de monotonicitat), el qual pot provocar que un candidat perdi
    l'elecció de forma paradoxal pel simple fet d'obtenir més suport en primera preferència per part d'un grup
    d'electors.
    Tampoc compleix necessàriament amb el criteri de Condorcet.
\end{itemize}

Davant d'aquestes alternatives, s'ha seleccionat el \textbf{mètode Schulze}.
Aquest algorisme realitza comparacions per parelles directes entre tots els candidats i construeix un graf de victòries
orientades.
Schulze és un mètode de Condorcet (sempre selecciona el guanyador de Condorcet si existeix) i compleix simultàniament
amb la monotonicitat, la independència de clons i la simetria de reversió, convertint-se en el mètode més robust i
equitatiu per al domini de Demochain.

\subsection{Inconvenients del Mètode Schulze}
\label{subsec:inconvenients-schulze}

Tot i ser excel·lent des del punt de vista de la justícia matemàtica, vàrem identificar quatre inconvenients principals
en emprar Schulze per a un sistema real:

\begin{enumerate}
    \item \textbf{Elevada complexitat computacional ($O(n^3)$)}: La determinació dels camins més forts en el graf de
    preferències per a resoldre els cicles (paradoxa de Condorcet) es basa en l'algorisme de Floyd-Warshall.
    El cost computacional creix de forma cúbica amb el nombre de candidats ($n$).
    Per mor d'això, calcular el recompte directament dins d'un contracte d'Algorand és completament inviable a causa
    dels límits de cost d'\textit{opcode} i temps d'execució de la màquina virtual, forçant a delegar el càlcul a
    l'aplicació client.

    \item \textbf{Baixa llegibilitat cognitiva per a la ciutadania}: La majority de votants està familiaritzada amb
    recomptes senzills (sumar vots directes).
    Entendre la lògica de camins de preferència més forts en un graf de duels Directed Acyclic Graph (DAG) i verificar
    que els resultats són correctes sense dependre d'un ordinador és gairebé impossible per al públic general.
    Això pot generar desconfiança o distanciament psicològic cap a la legitimitat de l'urna electrònica.

    \item \textbf{Manca d'escalabilitat del missatge de vot}: Cada elector ha d'ordenar les seves preferències.
    La quantitat d'informació necessària per a codificar el vot creix de forma quadràtica amb relació al nombre de
    candidats.
    Per a eleccions amb molts candidats, la \textit{payload} de la transacció i l'emmagatzematge en \textit{boxes} de la
    \textit{blockchain} creixen de forma ineficient, augmentant l'MBR a Algorand.

    \item \textbf{Complexitat en les regles de desempat}: Malgrat que l'algorisme està dissenyat per a definir clarament
    els camins de victòria, en escenaris reals es poden produir empats on els valors dels camins més forts de dos
    candidats siguin exactament idèntics.
    Resoldre aquests empats de manera justa exigeix implementar protocols deterministes addicionals preestablerts (com
    analitzar preferències de primer nivell o sortejos pseudoaleatoris basats en \textit{hash}es \textit{on-chain}),
    augmentant la dificultat del codi.
\end{enumerate}


\section{Llenguatge Ubic (DDD)}
\label{sec:llenguatge-ubic}

Per tal de garantir que tot l'equip de programació, els custodis i els auditors utilitzin el mateix significat de les
dades, s'ha adoptat una aproximació de \textbf{Domain-Driven Design (DDD)} definint un llenguatge comú per a tot el
projecte.
El domini de Demochain s'estructura a partir de la gènesi d'una \textbf{Organització}, la qual representa l'entitat
jurídica o el col·lectiu registrat al sistema per un \textbf{Organitzador} que és l'administrador amb drets d'escriptura
per gestionar les adreces del \textbf{Cens} d'aquella entitat.
Cada membre autoritzat d'aquest cens rep el rol de \textbf{Votant} i té dret a participar en les consultes de les
propostes.
La \textbf{Proposta} constitueix la consulta electoral \textit{on-chain}, que s'inicia sota una \textbf{Fase
d'Aprovació} on es requereix un quòrum de dues terceres parts favorables per transitar cap a la \textbf{Fase d'Elecció}
activa, durant la qual s'admet l'enviament de paperetes.

La manifestació del sufragi del votant s'anomena \textbf{Papereta}, la qual conté el missatge digital signat amb les
seves preferències, ja sigui com un vot de si/no o com un \textit{array} d'índexs ordinals.
En finalitzar el procés, es duu a terme el \textbf{Recompte}, que en el disseny teòric representa l'agregació
homomòrfica multiplicativa \textit{on-chain} de les paperetes xifrades sota una clau de desxifratge de llindar
controlada per diversos \textbf{Custodis o \textit{Trustees}}.
Cada custodi aporta la seva \textbf{\textit{Share} de Desxifratge} de forma independent sota un esquema de
\textbf{Consens \textit{K-of-N}}.
Per prevenir el doble vot de manera anònima, es preveu l'ús d'un \textbf{Nullifier o Anul·lador} determinista.
En l'\ac{MVP} simplificat actual, las paperetes es desen en clar i el recompte es calcula \textit{client-side},
mantenint l'estructura terminològica com a model per a futures iteracions.
Tot el cicle es tanca amb l'\textbf{Ancoratge}, que és l'acció determinista de publicar el \textit{hash} de resultats
finals d'Algorand a la xarxa d'Ethereum Sepolia per certificar-ne la integritat a llarg termini.


\section{Model de Dades}
\label{sec:model-dades}

El disseny de les dades del projecte s'ha estructurat de manera que encaixi amb les limitacions d'emmagatzematge
transaccional d'ambdues xarxes.

\subsection{Estructures d'Algorand (\texttt{ARC-4 Structs} en \texttt{algopy})}
\label{subsec:estructs-algorand}

Al contracte intel·ligent d'Algorand escrit en Python, s'utilitzen estructures de dades \ac{ARC}-4 que es codifiquen de
forma compacta en l'estat de les \textit{boxes}.
Les estructures estan definides a \texttt{contract.py}:

L'estructura \texttt{Organization} s'encarrega d'emmagatzemar les metadades generals de cada col·lectiu registrat al
sistema.
Conté un identificador numèric de l'entitat, el nom textual de l'entitat, una breu descripció, l'adreça de
\textit{wallet} del seu organitzador i el nombre total de membres que formen el cens.
Aquestes estructures es classifiquen i es desen en un mapa de caixes de dades persistents associades a l'estat del
contracte sota el prefix de clau d'organització:

\begin{pyverbatim}
    class Organization(arc4.Struct):
    org\_id: arc4.UInt64
    name: arc4.String
    description: arc4.String
    organizer: arc4.Address
    member\_count: arc4.UInt32
\end{pyverbatim}

Per altra banda, l'estructura \texttt{Proposal} recull el cicle de vida de cada elecció.
Emmagatzema l'identificador de l'organització a la qual pertany, el títol, la descripció del sufragi, un \textit{array}
dinàmic de cadenes de text amb la llista de candidats o opcions, i les dates d'inici i de tancament de la finestra de
votació representades en format UNIX:

\begin{pyverbatim}
    class Proposal(arc4.Struct):
    org\_id: arc4.UInt64
    title: arc4.String
    description: arc4.String
    options: arc4.DynamicArray[arc4.String]
    starting\_date: arc4.UInt64
    ending\_date: arc4.UInt64
\end{pyverbatim}

Per dur el registre dels sufragis d'aprovació i controlar la transició cap a la fase electoral s'empra l'estructura
\texttt{ApprovalTally}, que conté el nombre total de vots emesos a favor de la proposta i el nombre acumulat d'electors
participants en la fase d'aprovació.
Finalment, es defineixen les estructures \texttt{BallotId} i \texttt{CensusId} per actuar com a claus compostes
\textit{on-chain}, les quals s'encarreguen d'enllaçar de forma unívoca l'adreça del votant amb la proposta de votació
preferencial o amb la llista del cens de l'organització:

\begin{pyverbatim}
    class ApprovalTally(arc4.Struct):
    votes\_for: arc4.UInt32
    total\_votes: arc4.UInt32

    class BallotId(arc4.Struct):
    sender: arc4.Address
    proposal\_id: arc4.UInt64

    class CensusId(arc4.Struct):
    org\_id: arc4.UInt64
    member: arc4.Address
\end{pyverbatim}

En l'\ac{MVP} desenvolupat, el cens s'emmagatzema com un \texttt{BoxMap} que fa \textit{mapping} de cadascun dels
\texttt{CensusId} amb un valor booleà, i les paperetes preferencials en clar es desen directament com un \texttt{BoxMap}
que enllaça cadascun dels \texttt{BallotId} amb un \textit{array} dinàmic d'enters sense signe
(\texttt{arc4.DynamicArray[arc4.UInt8]}) que reflecteix el rànquing.

\subsection{Estructures de Dades d'Ethereum (\texttt{NotaryContract.sol})}
\label{subsec:estructs-ethereum}

Al contracte intel·ligent escrit en Solidity d'Ethereum, s'emmagatzema l'estat del consorci electoral per a cada sufragi
finalitzat mitjançant l'estructura \texttt{Election}.
Aquesta estructura recull el llindar \texttt{thresholdK} necessari per segellar el resultat, un \textit{array} dinàmic
amb les adreces públiques de les universitats certificadores del consorci, i mapes de correspondència per validar quins
nodes estan autoritzats a firmar, quines submissions s'han registrat des de les seves adreces de Sepolia i l'acumulació
de vots que té cada variant de \textit{hash}.
L'estructura finalitza amb dos valors booleans que indiquen si l'elecció es troba actualment oberta a l'ancoratge o si
ja ha assolit el quòrum i es troba oficialment ancorada a la xarxa:

\begin{verbatim}
struct Election {
    uint256 thresholdK;
    address[] authorizedNodes;
    mapping(address => bool) isAuthorized;
    mapping(address => bytes32) submissions;
    mapping(bytes32 => uint256) \textit{hash}Count;
    bool open;
    bool anchored;
}
\end{verbatim}

\subsection{Procés de Derivació del \textit{hash} Canònic}
\label{subsec:derivacio-hash}

El \textit{hash} que s'envia a Ethereum Sepolia representa el compromís d'integritat del resultat electoral.
Per calcular-lo de manera idèntica a tots os nodes validadors s'aplica el següent procés determinista:

\begin{enumerate}
    \item S'escanegen totes les \textit{boxes} d'Algorand com a llista de paperetes amb prefix \texttt{eb\_} vinculades
    a la proposta.

    \item Es classifiquen les paperetes ordenant les adreces dels electors de forma lexicogràfica.
    Aquest pas és vital, ja que l'ordre físic de rebuda de transaccions depèn de la latència de la xarxa i canviaria el
    format en bytes de cada node.

    \item Se serialitza la informació en una cadena de text tipus \ac{JSON} canònica (claus ordenades de la A a la Z,
    sense cap espai en blanc, tabulacions o salts de línia):
    \begin{verbatim}
        {
            "ballots": [
                {
                    "preference": [2, 0, 1],
                    "voter":"ADDR_A"
                },
                {
                    "preference": [0, 2, 1],
                    "voter":"ADDR_B"
                }
            ],
            "options": ["Opció A", "Opció B", "Opció C"],
            "proposal_id": 1,
            "title": "Rector 2026"
        }
    \end{verbatim}

    \item Es calcula el \textit{hash} \ac{SHA}-256 d'aquesta representació en bytes, produint un valor únic hexadecimal
    de 32 bytes (\texttt{bytes32}).
\end{enumerate}
